\documentclass[11pt]{article}
\usepackage{mathblatt}
% gesetzt von werkzeuge/setzer.py; Übersicht nach bau/bauregeln.md „Übersicht und Serie“.
\newcommand{\kennung}{SLM}
\newcommand{\bereich}[1]{\par\addvspace{14pt}\noindent{\large\bfseries #1}\par\nobreak\vspace{4pt}}
\newcommand{\bl}[1]{\par\noindent #1\par\vspace{3pt}}
\newcommand{\blrand}[1]{\par\noindent{\color{mbgrau}#1}\par\vspace{3pt}}
\newcommand{\blhier}[1]{\par\noindent\llap{$\blacktriangleright$\hspace{0.5em}}\textbf{#1}\par\vspace{3pt}}
\begin{document}
\pfheftstil
\fancyfoot[C]{{\footnotesize\color{mbgrau}\kennung\hfill\thepage}}
\pfheftkopf{Potenz- und Exponentialfunktionen, Wachstum und Zerfall}{Klasse 10}

\bereich{Vorher}
\bl{Prozentrechnung: Wachstumsfaktor}
\bl{Lineare Funktionen}

\bereich{Wachstum und Zerfall}
\bl{Lineares und exponentielles Wachstum unterscheiden}
\blhier{Wachstumsfaktor und Wachstumstabelle}
\bl{Exponentialfunktion aufstellen und auswerten}
\bl{Verdopplungs- und Halbwertszeit}

\bereich{Weiter}
\bl{Potenzfunktionen}
\end{document}
